\chapter{Conclusions}
\label{ch:conclusions}
\doublespacing
\vspace{-2.5cm}

To provide a closing remark on the work hereby presented, it is worth to summarize the main contributions and results obtained, also to highlight possible improvements to consider and future directions to explore. This work addressed the context of distributed autonomous systems with the purpose of designing an architecture supporting research activities and workflows involving them. The importance of such an architecture, as well as its basic requirements, have been illustrated in \Cref{ch:prb}. In \Cref{sec:prb_related_work}, it has been recognized how one of the major drawbacks of existing solutions is the intent to provide a complete implementation from the ground up, instead of relying on existing tools and technologies that help to solve at least a part of the issues at hand. The relevant technologies have been introduced in \Cref{ch:tools}, highlighting for each the main features that make them suitable for the design of the architecture. The Distributed Unified Architecture has been presented in \Cref{ch:dua} as a framework born from the integration of existing technologies with specific, novel solutions. The discussion addressed each of the basic requirements with a DUA feature, showing how the proposed architecture provides a common base environment for the development of software modules for autonomous systems, on both PCs and dedicated hardware.

The proposed solution has been tested on a variety of hardware platforms, whose interoperability has also been shown in the use cases considered in \Cref{ch:robots}. A possible improvement in this regard is certainly the development of new base units (see \Cref{sec:dua_foundation}): since the proposed approach worked on the set of hardware platforms discussed so far, it is worth expanding this set including new solutions and refining the existing ones; a first step in this regard would be, \emph{e.g.}, the development of a \texttt{jetson6} base unit supporting the most recent JetPack 6 SDK for the Nvidia Jetson platforms. Another possible improvement concerns the reliance on Docker (see \Cref{sec:docker,sec:dua_foundation}) even on SBCs and similar platforms. While the container-based approach of \texttt{dua\_foundation} clearly helps to abstract the hardware layer in a repeatable and reliable fashion, it adds a degree of complexity that is preferable to avoid in embedded systems. To this end, a restructuring of the \texttt{dua\_foundation} module could be considered, continuing to provide the development base units in the form of Docker images but developing the ones targeted to embedded systems as classic system images using, \emph{e.g.}, crosstool-NG \cite{morin_crosstool-ng_nodate} and Buildroot \cite{korsgaard_buildroot_nodate} to build the toolchains and the root filesystems, respectively. Another aspect to consider is that, as discussed in \Cref{sec:dua_remarks}, the Development (D) requirement is still not completely satisfied by the proposed architecture. Two areas for improvement are identified in this context: the addition and configuration of cross-compilation toolchains and debuggers in the development base units, and the integration of support for Continuous-Integration/Continuous-Deployment (CI/CD) pipelines in the \texttt{dua\_template}, to support the most modern automated development workflows.

The many benefits of such an architecture in the autonomous robotics field have been discussed in \Cref{ch:robots}, with some results highlighted in \Cref{sec:robots_results}. To this end, the improvement of the development tools that DUA offers to ROS 2 and robotic software developers is certainly a path to consider. Among the many possible outcomes of such an effort, some notable ones include the development of a version of the \texttt{params\_manager} module supporting Python packages, and a mathematical library for the management of the state of rigid bodies in 3D space offering APIs to interface directly with ROS 2 communication primitives. More in general, getting the community more involved in the DUA project would certainly be beneficial, for both the project itself and the ROS 2 ecosystem: some of the features that \texttt{dua\_utils} modules currently offer could be integrated into the ROS 2 API, and the development of new modules could be driven by the broader needs of the community.

The other application scenario considered in \Cref{ch:tokamaks,ch:allocation}, that is, that of magnetic confinement nuclear fusion, also provides some notable results. In \Cref{ch:tokamaks}, the application of the DDS middleware to the distributed control system of the TCV tokamak clarified the importance of these technologies even in critical research environments as that of the SCD distributed control system. These solutions also solved critical constraints affecting the expandability of the research infrastructure, paving the way for the integration of new real-time diagnostic systems and control algorithms. The benchmarking results shown in \Cref{sec:ddsdatasource_benchmark} confirmed that the intergation of the DDSDataSource in the SCD architecture, to replace the previous reflective memory network, is an update to be performed. Furthermore, the results concerning the dynamic allocation of EF coil currents in \Cref{ch:allocation} also deserve further investigation. The methodological solutions developed in \Cref{sec:all_design} cover a wide variety of scenarios but also uncover further questions to be researched. First, the study presented in \Cref{sec:all_design_ssinv} testifies that the dynamic input allocation problem for nonlinear systems is meaningful and may also be solved by framing it in the context of optimal control; however, further developments are necessary to understand to what extent this approach is conclusive. The dynamic annihilator presented in \Cref{sec:all_design_fullinv} and introduced for the first time in \cite{cristofaro_output_2014} broke the tight coupling involving full output invisibility and state-space properties since \cite{zaccarian_dynamic_2009} but also shed light of more questions; more research is required to investigate how input-output characterization of the invisibility property may be reflected on state-space properties of the plant. Furthermore, the sampled-data case has been addressed for the first time in the literature as reported in \Cref{sec:all_design_sd} and many questions remain open, such as the proper definition of the inter-sample output ripple due to the sampled input allocation action. The implementation of a sampled-data dynamic allocator may also be performed on the TCV control system, extending the developments of \Cref{sec:tcvall} and providing, hopefully, a first practical proof of its effectiveness. Finally, the experimental results provided in \Cref{sec:shots}, obtained with the dynamic EF coil current allocator designed in \Cref{sec:tcvall_design}, constitute the first practical example of the effectiveness of a dynamic control allocator designed for full output invisibility, translating the theoretical results of \Cref{sec:sota_allocation,sec:all_design_fullinv} into a real-world application involving a sophisticated plant, such as the TCV tokamak, whose technical constraints make it benefit from the proposed solution. Another way to provide more experimental results corroborating the methodological approaches mentioned above is to test the proposed dynamic allocators on a wider variety of plasma shapes and optimization objectives.
